\section{Proof of the unequal-product bound and the local-completion limit}
\label{upq:details}

\subsection{The original expansion before absolute Cauchy}
Let $\mathcal P_U$ consist of the ordered pairs of good squarefree ideals
$(a,b)$ with
\[
 (a,b)=1,\qquad Na,Nb>U,\qquad D\le Na\,Nb\le2D.
\]
The hyperbolic and coprimality conditions are part of this set. The cutoff
window implies $Na,Nb\asymp\sqrt D$, with constants independent of $U$.
Put
\[
 c_{a,b}=\mu(a)\mu(b)\nu(ab)W(Nab/D),\qquad
 \mathcal C(u)=\sum_{(a,b)\in\mathcal P_U}c_{a,b}x_a(u)x_b(u),
 \qquad
 \mathcal M(u)=\sum_n\lambda_nx_n(u).
\]
Because $ab$ is squarefree, $\mu(a)\mu(b)=\mu(ab)$ and
$x_a(u)x_b(u)=x_{ab}(u)$ including every zero. Hence
$\mathcal C(u)=\sum_mr_m\lambda_mx_m(u)$ is exactly the original
$C_U$ polynomial. In particular, the exact expansions are
\begin{align}
 \frac1D\sum_{u\in B_A}|\mathcal C(u)|^2
 &=\frac1D\sum_{\substack{(a,b)\in\mathcal P_U\\
                          (c,d)\in\mathcal P_U}}
 c_{a,b}\overline{c_{c,d}}
 \sum_{u\in B_A}x_a(u)x_b(u)
       \overline{x_c(u)x_d(u)},                         \label{upq:four}\\
 \frac1D\sum_{u\in B_A}\mathcal C(u)\overline{\mathcal M(u)}
 &=\frac1D\sum_{\substack{(a,b)\in\mathcal P_U\\n}}
 c_{a,b}\overline{\lambda_n}
 \sum_{u\in B_A}x_a(u)x_b(u)\overline{x_n(u)}.           \label{upq:mixed}
\end{align}
Grouping the two squarefree products gives the sectors in the core.
No internal mask or common profile has been discarded before this
grouping. An unmasked product of four transforms would change the
support; signed coefficients do not give pointwise domination merely
by allowing additional factorizations.

\subsection{A legal annular norm-pair count}
For $m\ne n$ in the annulus, write uniquely
\[
 h=(m,n),\qquad m=hr,\qquad n=ht.
\]
Squarefreeness makes $h,r,t$ pairwise coprime and squarefree. The conditions
$D\le Nm,Nn\le2D$ imply
\[
 \frac12\le\frac{Nr}{Nt}\le2,\qquad
 Nh\max(Nr,Nt)\le2D.
\]
If $Nr\,Nt\le Y$ and $V=\max(Nr,Nt)$, then
\[
 V^2\le2Y,\qquad V\le2D.
\]
On a dyad $V_0\le V<2V_0$, elementary ideal counting gives
$O_K(V_0^2)$ pairs $(r,t)$. For each such pair there are separately
$O_K(D/V_0)$ choices of $h$. This last estimate remains valid on every
nonempty terminal range, since $V_0\le2D$; an upper norm limit below one
contains no ideals. Enlarging the positive count to arbitrary ideals
can only increase it. Thus this dyad contributes $O_K(DV_0)$, and
\[
 \sum_{\substack{V_0\ {\rm dyadic}\\V_0\le\sqrt{2Y}}}
       DV_0\ll_K D\sqrt Y.
\]
We have proved
\begin{equation}\label{upq:paircount}
 \#\{(m,n):m\ne n,\ D\le Nm,Nn\le2D,\
             N(mn/(m,n)^2)\le Y\}\ll_K D\sqrt Y.
\end{equation}
The sum is geometric, so no logarithm is needed. The independent
$O_K(D^2)$ count can be used when $Y$ exceeds the possible conductor
range. This is an unequal-product annular count, not the earlier
coefficient-diagonal multiplicative-energy estimate.

For every $\delta>0$, on squarefree ideals
\[
 r_U(m)\le2^{\omega(m)}\ll_{K,\delta}(Nm)^\delta.
\]
To see the last inequality, choose a finite prime set so that
$2\le(Np)^\delta$ outside it; the remaining finite product is a constant.
Choose $\delta$ small enough to absorb both factor multiplicities into
$D^\epsilon$. The exact row kernel satisfies
$|K_A(m,n)|\le\#B_A\ll_K R_0$, retaining all literal row restrictions.
Multiplying~\eqref{upq:paircount} by this bound, the actual multiplicities,
and $\|W\|_\infty^2$, then dividing by $D$, proves~\eqref{upq:sc}.
For the mixed expression only one multiplicity $r_U(m)$ occurs.
The estimate bounds its complex absolute value, hence also its real part.
Its constants are uniform in the high cutoff window and use no derivative
of the possibly complex common $W$.

\subsection{The whole saturation face and its neighborhood}
Write $C=H^{1/3}$. On~\eqref{upq:face},
\[
 R_0=CT,\qquad P=\max(T,C/T),\qquad
 R_0P=\max(CT^2,C^2)\le C^3=H.
\]
This proves the face assertion at $Y=P^2$ for every $1\le T\le C$;
the endpoint $A_3=1$ is not singled out.

For completeness, let $x_j=\log A_j$, $r_0=\sum_jx_j$,
$c_0=\sum_jjx_j\le\log H$. The accepted exact slack identity is
\[
 \log H-3\{r_0-\max(x_1,x_2)\}
 =\log H-c_0+
 \begin{cases}
 x_1-x_2+x_4+2x_5+3x_6,&x_1\ge x_2,\\
 2(x_2-x_1)+x_4+2x_5+3x_6,&x_2\ge x_1.
 \end{cases}
\]
All terms are nonnegative. If $R_0/M\ge H^{1/3}D^{-\eta}$, then
$A_1/A_2\le D^{3\eta}$ in the first branch; in the other branch that ratio
is already at most one. For every $j\ge2$,
$R_0A_j\le\prod_iA_i^i\le H$, whereas
\[
 \frac{R_0A_1}{\prod_iA_i^i}
 =\frac{A_1}{A_2A_3^2A_4^3A_5^4A_6^5}
 \le\frac{A_1}{A_2}.
\]
Consequently $R_0P\le H\max(1,A_1/A_2)\le HD^{3\eta}$.
This proves the claimed neighboring bound with the complete $A_3$
direction, both wedges and the $A_4,A_5,A_6$ slack retained.
No principal residue has been redefined.

\subsection{Original versus arbitrary four-transform conductors}
At a good prime $p$ let
\[
 (\alpha,\beta,\gamma,\delta)
 =(\mathbf1_{p\mid a},\mathbf1_{p\mid b},
   \mathbf1_{p\mid c},\mathbf1_{p\mid d}),\qquad
 e=\alpha+\beta-\gamma-\delta.
\]
For $p\nmid u$ the local product in~\eqref{upq:four} is
$\chi_p(u)^e$. For $p\mid u$ it is zero whenever any membership is
nonzero, including when $e=0$. For four arbitrary individually squarefree
factors, $e\in\{0,\pm1,\pm2\}$. The nonzero exponents have nontrivial tame
local characters, so their conductor exponent is one. At $e=\pm1$
the order is six, and at $e=\pm2$ it is three.

On the original support, $\alpha+\beta,\gamma+\delta\in\{0,1\}$.
Thus $e=\pm2$ cannot occur and
\[
 q=\prod_{\substack{p\mid L\\e_p\ne0}}p
   =\operatorname{supp}(ab)\mathbin{\triangle}
       \operatorname{supp}(cd),\qquad
 L=\operatorname{rad}(abcd).
\]
Here the support expression denotes its squarefree ideal product.
With $h=(ab,cd)$, $L=hq$ and the exact kernel factor is
\begin{equation}\label{upq:rawpair}
 x_{ab}(u)\overline{x_{cd}(u)}
 =\mathbf1_{(u,h)=1}\theta_{ab/h,\,cd/h}(u).
\end{equation}
The primitive character on the right already vanishes at primes of $q$;
the explicit indicator supplies every canceled common-prime zero.
Split conjugate prime ideals are distinct ideal factors even when their
rational norms agree. The defining fixed $S$ data are unchanged.

For an unrestricted membership $(1,1,0,0)$, $e=2$ is cubic, although the
parity symmetric difference of the four supports is empty.
For $(1,1,1,0)$, $e=1$, although $p$ lies in both nonsquarefree product
supports. Neither relaxed convention is the original two-squarefree-product
formula. In~\eqref{upq:mixed} the original $(a,b)=1$ similarly restricts
$e=\mathbf1_{p\mid a}+\mathbf1_{p\mid b}-\mathbf1_{p\mid n}$
to $0,\pm1$.

Fix the other row factors and pivot in direction $j$. The local exponent
becomes $je$ modulo six. The complete order/conductor table is
\[
 \begin{array}{c|c|c}
 j&e=\pm1&e=\pm2\\ \hline
 1,5&\text{order }6,\ \text{conductor }p&
       \text{order }3,\ \text{conductor }p\\
 2,4&\text{order }3,\ \text{conductor }p&
       \text{order }3,\ \text{conductor }p\\
 3&\text{order }2,\ \text{conductor }p&
       \text{trivial,\ conductor }1\\
 6&\text{trivial,\ conductor }1&
       \text{trivial,\ conductor }1
 \end{array}
\]
Every cell retains a raw zero at $p\mid a_j$ if $p\mid L$.
For general four factors the primitive support is $Q_j$ and the additional
deleted ideal is $Z_j=L/Q_j$ as in the core. If a fixed other row factor
meets $L$, the whole original summand is zero before the pivot is summed.
For the original unequal products, $Q_j=q$ for $j=1,\ldots,5$.
In particular the full $a_3$ direction is quadratic with conductor $q$,
not primitive-principal. Its smaller character order does not shorten its
conductor. In a relaxed $\pm2$ channel the $j=3$ component is
primitive-principal, but still requires its union deletion.

\subsection{Complete local sums and the full deletion union}
Let $\chi$ be the order-six character on a good residue field
$\mathbb F_Q$ and $\rho$ a nontrivial additive character. Define
\[
 G_k(t)=\sum_{z\in\mathbb F_Q^\times}\chi(z)^k\rho(tz).
\]
For $k\not\equiv0\pmod6$, character orthogonality gives $G_k(0)=0$.
For $t\ne0$, substitution gives
$G_k(t)=\overline{\chi(t)^k}G_k(1)$. Expanding its square and putting
$x=ay$ yields
\[
 |G_k(t)|^2
 =\sum_{a\ne0}\chi(a)^k
       \sum_{y\ne0}\rho\{t(a-1)y\}
 =(Q-1)-\sum_{a\ne0,1}\chi(a)^k=Q.
\]
For $k=0$ the omitted zero gives $G_0(0)=Q-1$ and $G_0(t)=-1$
when $t\ne0$. This proof applies both to split prime fields and to inert
quadratic residue fields.

The Chinese remainder theorem expresses the full period-$L$ transform
as these local factors. Its additive identifications may introduce unit
twists, which change phases but not magnitudes. It is zero if any local
frequency on $Q_j$ vanishes; otherwise its magnitude is
\begin{equation}\label{upq:gaussunion}
 (NQ_j)^{1/2}
 \prod_{p\mid Z_j}
 \left|-1+Np\,\mathbf1_{\text{local frequency at }p\text{ is }0}\right|.
\end{equation}
Thus the full transform is not uniformly bounded by $\sqrt{NQ_j}$
when $Z_j\ne1$. A canceled prime supplies a Ramanujan factor.
The zero full frequency vanishes for $Q_j>1$, but a short original
row window has nonzero-frequency contributions.

For a fixed pivot, with its remaining support and restrictions denoted
by $F$, exact deletion gives
\begin{align}
 \mathbf1_{(r,Z_j)=1}
 &=\sum_{\substack{d\mid Z_j\\d\mid r}}\mu(d),\nonumber\\
 \sum_r\theta_j(r)F(r)\mathbf1_{(r,Z_j)=1}
 &=\sum_{d\mid Z_j}\mu(d)\theta_j(d)
       \sum_z\theta_j(z)F(dz).                         \label{upq:lift}
\end{align}
Use fixed primary generators for ideal multiplication and retain all
literal unit choices. The scaled $F(dz)$ includes the actual ray class,
window and row masks; this identity asserts no independence.

For a squarefree pivot let $E$ be the radical of $Z_j$ and its other
fixed excluded primes, including the relevant other $a_i$ with $i\le5$.
Remove primes of $Q_j$, where $\theta_j$ already vanishes. Pointwise,
\[
 \mu(r)^2\mathbf1_{(r,E)=1}
 =\sum_{\substack{h^2\mid r\\(h,E)=1}}\mu(h)
       \sum_{\substack{d\mid E\\d\mid r}}\mu(d).
\]
If $(r,E)>1$, the second sum is zero; otherwise the first is exactly
the squarefree indicator. Since $(d,h)=1$, the pivot sum becomes
\begin{equation}\label{upq:sflift}
 \sum_{d\mid E}\mu(d)\theta_j(d)
 \sum_{(h,E)=1}\mu(h)\theta_j(h)^2
       \sum_z\theta_j(z)F(dh^2z).
\end{equation}
Terms with $Nh>\sqrt{O(A_j/Nd)}$ vanish. The factor $dh^2$ is valid
because of this coprimality; overlapping divisibility conditions have
not been multiplied incorrectly. An arbitrary $a_6$ pivot has no
squarefree expansion.

The divisor triangle costs at most
$\tau(E)\ll_{S,\epsilon}(HD)^\epsilon$.
The elementary length bound in~\eqref{upq:sflift} is
$O(A_j/(Nd\,Nh^2))$ on its nonempty ranges. The reciprocal
norm-square sum over $h$ converges by ideal counting. Restoring the
deletion and squarefree restrictions therefore returns at best
$O((HD)^\epsilon A_j)$ by this calculation; it does not produce a power
saving.

\subsection{An exact completion identity and the favorable smooth trial}
For the literal sharp window one can periodize its remaining weight
$F$ modulo $L$. Finite Fourier inversion is exact:
\[
 \sum_r F(r)f_L(r)
 =\frac1{NL}\sum_t\widehat F_L(t)G_L(t),
 \qquad
 f_L(r)=\theta_j(r)\mathbf1_{(r,Z_j)=1},
\]
with mutually dual additive signs. Here $\widehat F_L$ includes every
actual restriction not already in $f_L$.
Although $G_L(0)=0$ for $Q_j>1$, the nonzero frequencies and
their factors~\eqref{upq:gaussunion} still require a bound.
The completed zero mean alone is not such a bound.

To test the best elementary local scale, consider instead an unmasked
smooth primitive sum
\[
 S_\theta(Y)=\sum_{z\in\mathbb Z[\omega]}
       \theta(z)V(z/\sqrt Y),\qquad Y\ge1,
\]
where $V$ is a fixed compactly supported smooth function on the
Eisenstein plane and $\theta$ has squarefree conductor $Q>1$.
A finite fixed-modulus decomposition retains the specified primary/ray
restrictions. Since the Eisenstein ring is a PID, the lattice of $Q$
is a rotated and scaled copy of the fixed lattice, with scale
$\sqrt{NQ}$. Periodization and Poisson inversion, the complete Gauss
bound, and rapid Fourier decay give, for any fixed $A>2$,
\[
 |S_\theta(Y)|
 \ll_{K,A,V}\frac{Y}{\sqrt{NQ}}
       \sum_{\substack{\xi\in\mathcal L^*\\\xi\ne0}}
       \left(1+\sqrt{\frac{Y}{NQ}}|\xi|\right)^{-A}.
\]
The fixed dual lattice $\mathcal L^*$ absorbs the Fourier normalization;
rotations do not affect the stated bound. If $Y\le NQ$, the lattice sum
is $O_A(NQ/Y)$ by planar counting. If $Y\ge NQ$, absence of the zero
frequency gives $O_A((NQ/Y)^{A/2})$. Together with the elementary
length estimate this proves the favorable bound
\begin{equation}\label{upq:smoothpoisson}
 |S_\theta(Y)|\ll_{K,V}\min(Y,\sqrt{NQ}).
\end{equation}
Bounded terminal norm scales are absorbed in constants.

This is a smooth local trial, not a replacement of the original sharp
cutoff in~\eqref{upq:four}--\eqref{upq:mixed}.
Its constant depends on fixed smooth seminorms of $V$ and need not remain
bounded when a sharp boundary is approximated. For the original sharp
masked kernel the unconditional bound obtained here is still
$|K_A(m,n)|\ll_KR_0$.

Even the favorable trial gives no improvement for the remaining
$Nq>P^2$ sector. On the original support $Q_j=q$ for $j\le5$, and
\[
 \sqrt{Nq}>P\ge A_j\ge
       \frac{A_j}{Nd\,Nh^2}.
\]
Thus every scaled minimum in~\eqref{upq:smoothpoisson} takes its length
term when inserted into~\eqref{upq:lift} or~\eqref{upq:sflift};
the remaining masks must still be accounted for. Per-conductor
absolute summation loses the signed column information and returns
$R_0$ up to divisor losses. In the coprime-product subsector $(m,n)=1$,
$Nq=Nm\,Nn\in[D^2,4D^2]$. At $H=D^{3/4}$ on the complete face,
$A_1,A_2,A_3\le D^{1/4}$ but $\sqrt{Nq}\ge D$.
Neither the sextic, cubic nor quadratic local order changes that
conductor comparison.

\subsection{The failed inference and the remaining signed sum}
The accepted balanced-factor count gives
$\#\mathcal P_U=\sum_mr_U(m)\ll_KD$, and there are $O_K(D)$
single columns. Using only $|K_A|\ll R_0$ in the absolute off-diagonal
expansions therefore costs $O(\|W\|_\infty^2DR_0)$ for each expression.
At $B_X$ this is $DX^2$, whereas the inherited powered-BGL middle
term is $DX$. On the full face with $C=D^{1/4}$,
\[
 \begin{array}{c|c}
 \text{calculation}&\text{bound apart from constants and }\epsilon\\ \hline
 \text{fully absolute unequal products}&DCT\\
 \text{inherited middle monomial}&DC=D^{5/4}\\
 \text{new }Nq\le P^2\text{ sector}&
       \max(CT^2,C^2)\le D^{3/4}.
 \end{array}
\]
These are upper estimates, not lower bounds for the actual signed sums.

The unsupported step would be to infer short masked cancellation from
the complete mean zero. Its nonzero additive frequencies have not been
estimated at a useful scale, canceled-prime Ramanujan factors cannot be
omitted, and the exact sharp boundary cannot be silently smoothed with
a uniform constant. This is a failure of the displayed pointwise
completion calculation, not a disproof of the original signed target.

The surviving target is the original sum
\[
 \frac1D\sum_{\substack{m\ne n\\N(mn/(m,n)^2)>P^2}}
 r_mr_n\lambda_m\overline{\lambda_n}K_A(m,n),
\]
jointly with the expression having only $r_m$.
Keep the factor realization of each $r_m$, its full hyperbolic support,
and~\eqref{upq:rawpair} on every actual row block, including the full
face and neighboring region. A useful bound must combine the signed
column correlations with the nonzero dual frequencies before the
per-conductor absolute estimate. No such long-conductor bound is
proved here. Principal centering, the independently primitive
$\psi_u,\psi_{u,2}$, and all deleted sets retain their previous meanings.

The standard-library diagnostic in
\path{planning/unequal-product-20261105/attempt/}
checks exact finite split/inert residue arithmetic, Gauss/Ramanujan
identities, the original fourth and mixed expansions, union deletions,
squarefree lifts and the full face. Its common complex norm profiles
are realizable by smooth disjoint bumps. These finite controls do not
establish an asymptotic saving. No complete $BT_\sigma$ bound, improved
full exponent, new Lean theorem, zero-location result or GRH conclusion
is asserted.
